\FloatBarrier
\section{Rotating Frame}
\label{app:spin-driving}
\label{sec:calibration-spin-drive}
\label{sec:drive-confined-spin}
\label{sec:drive-rotating-frame}

Section~\ref{sec:transverse-spin-control} uses an effective
rotating-frame Hamiltonian to describe transverse spin driving.
Keeping the rapidly oscillating term during the transformation
makes explicit the approximation behind
Eq.~\eqref{eq:esr-propagator} and fixes the phase convention used
for the rotation axes.

Use the spin doublet of Section~\ref{sec:platform-encoding}, with
$\ket0=\ket{\downarrow}$ and $\ket1=\ket{\uparrow}$.
This ordering exchanges the two physical-spin basis states,
reversing the signs of the $z$ and $y$ Pauli operators while
leaving the $x$ operator unchanged:
\[
 \sigma_x=X,\qquad \sigma_y=-Y,\qquad \sigma_z=-Z.
\]
Allow the real drive coefficient in Eq.~\eqref{eq:esr-lab} to
have a slowly varying envelope $\Omega(t)$, keeping $\phi$ fixed.
In the logical basis, its Hamiltonian is
\begin{equation}
 H_{\mathrm{lab,q}}(t)
 =-\frac{E_Z}{2}Z
   +\hbar\Omega(t)\cos(\omega t+\phi)X.
 \label{eq:calibration-electric-spin-drive}
\end{equation}
Here $E_Z$ is the static splitting, $\omega$ the carrier angular
frequency, and $\Omega(t)$ the transverse drive coefficient in
angular-frequency units.

The drive in Eq.~\eqref{eq:calibration-electric-spin-drive}
oscillates at $\omega$, whereas near resonance the spin's motion
relative to that carrier is slow. To describe the latter motion,
write the laboratory state as
\[
 \ket{\psi_{\mathrm{lab}}(t)}
 =U_{\mathrm{rf}}(t)\ket{\psi_{\mathrm{rot}}(t)},\qquad
 U_{\mathrm{rf}}(t)=e^{i\omega tZ/2}.
\]
The rotating state evolves with
\begin{equation}
 H_{\mathrm{rot}}(t)
 =U_{\mathrm{rf}}^\dagger H_{\mathrm{lab,q}}(t)U_{\mathrm{rf}}
  -i\hbar U_{\mathrm{rf}}^\dagger\dot U_{\mathrm{rf}}.
 \label{eq:calibration-rotating-frame-transform}
\end{equation}
The second term contributes $+\hbar\omega Z/2$, reducing the
longitudinal splitting to the difference between the Zeeman and
carrier frequencies. The transverse operator becomes
\[
 U_{\mathrm{rf}}^\dagger XU_{\mathrm{rf}}
 =X\cos\omega t+Y\sin\omega t.
\]
Multiplying this expression by the carrier separates the drive
into a stationary part and a part oscillating at twice the
carrier frequency:
\[
 \begin{aligned}
 &\cos(\omega t+\phi)
       (X\cos\omega t+Y\sin\omega t)\\
 &\qquad=\frac12(X\cos\phi-Y\sin\phi)
       +\frac12\bigl[X\cos(2\omega t+\phi)
                    +Y\sin(2\omega t+\phi)\bigr].
 \end{aligned}
\]

When $|\delta|,|\Omega(t)|\ll\omega$ and the envelope changes
little during one carrier
period, the $2\omega$ term is fast compared with the rotation of
interest. Dropping it is the rotating-wave approximation (RWA).
The remaining Hamiltonian is
\begin{equation}
 \begin{aligned}
 H_{\mathrm{rot,q}}(t)&=\frac{\hbar}{2}
 \left[-\delta Z+
  \Omega(t)(\cos\phi X-\sin\phi Y)\right],\\
 \delta&=E_Z/\hbar-\omega.
 \end{aligned}
 \label{eq:calibration-single-spin-hamiltonian}
\end{equation}
On resonance, a constant envelope therefore produces rotations
at angular frequency $|\Omega|$. This is the Rabi frequency
associated with the laboratory-frame coefficient defined above.

Equation~\eqref{eq:calibration-single-spin-hamiltonian} also
identifies the role of each control: the frequency fixes the
detuning, the phase fixes the transverse rotation axis, and the
envelope fixes the rotation strength. For positive $\Omega$,
$\phi=0,-\pi/2,\pi,+\pi/2$ select $+X,+Y,-X,-Y$, respectively.
Turning the envelope off leaves the waiting Hamiltonian
\[
 H_{\mathrm{wait}}=-\frac{\hbar\delta}{2}Z.
\]
These expressions use a constant reference frequency $\omega$
and neglect noise in its phase. With that reference held fixed,
a time-dependent detuning comes from a changing spin splitting
$E_Z(t)$. Population relaxation requires the open-system transition
rates described in Section~\ref{sec:measurement-relaxation};
it is not generated by these coherent two-state Hamiltonians.

To recover the constant-envelope result of the main text, set
$\Omega(t)=\Omega$ and substitute
$\sigma_z=-Z$, $\sigma_x=X$ and $\sigma_y=-Y$.
The transformation becomes $R(t)=e^{-i\omega t\sigma_z/2}$,
with $\ket{\psi(t)}=R(t)\ket{\psi_{\mathrm r}(t)}$.
For the laboratory Hamiltonian in Eq.~\eqref{eq:esr-lab},
keeping the $2\omega$ term gives the full expression
\begin{equation}
 \begin{aligned}
 H_{\mathrm r}
 &=R^\dagger HR-i\hbar R^\dagger\dot R\\
 &=\frac{\hbar\delta}{2}\sigma_z
   +\frac{\hbar\Omega}{2}
       (\cos\phi\,\sigma_x+\sin\phi\,\sigma_y)\\
 &\quad+\frac{\hbar\Omega}{2}
       \bigl[\cos(2\omega t+\phi)\sigma_x
             -\sin(2\omega t+\phi)\sigma_y\bigr],
 \qquad \delta=\omega_Z-\omega.
 \end{aligned}
 \label{eq:esr-rotating}
\end{equation}
Here $\omega_Z=E_Z/\hbar$. Omitting the last term gives precisely
the constant Hamiltonian used for the propagator in
Eq.~\eqref{eq:esr-propagator}.
